\section{Motivation}

Knowledge graph construction remains a critical bottleneck in computational drug discovery. Traditional approaches rely on domain-specific named entity recognition and relation extraction models trained on manually annotated corpora --- a process that is both expensive and brittle. A model trained to extract oncology relationships from clinical trial reports will not generalize to neuroscience literature, rare disease case studies, or immunology abstracts. Each new therapeutic area demands fresh annotation campaigns, new training runs, and careful validation against domain ontologies. For pharmaceutical researchers who routinely work across multiple therapeutic areas, this per-domain cost creates a prohibitive barrier to building the structured knowledge bases that modern drug discovery pipelines require.

In prior work approximately ten months ago, prompted by Microsoft Research's announcement of GraphRAG \cite{edge2024graphrag}, I explored whether large language models could circumvent this bottleneck entirely. Using GraphRAG's approach to query-focused summarization via graph construction from unstructured text, I built a latent knowledge graph from six drug discovery articles spanning multiple therapeutic areas~\cite{bhatt2025latentkg}. Using local embeddings via Ollama for indexing and OpenAI GPT-4o-mini via OpenRouter for query synthesis, the system extracted 3,224 entities and 2,242 relationships, organized into approximately 167 hierarchical community reports. The experiment proved the core concept: LLMs can extract structured knowledge from scientific literature without task-specific training or annotated corpora.

Despite these promising results, GraphRAG exposed critical limitations for scientific applications. While GraphRAG uses LLM extraction during its indexing phase, the resulting relationships lack domain-specific typing, ontological grounding, and evidence traceability --- the properties scientists require to trust and act on a knowledge graph. An edge between ``nivolumab'' and ``colitis'' does not indicate whether this represents an adverse effect, a treatment target, or merely a co-mention in the same passage --- yet this distinction is precisely what a drug safety scientist needs. Entities are generic NER output rather than identifiers grounded in biomedical ontologies; ``KRAS'' appears as a text span rather than a gene symbol mapped to HGNC, making downstream integration with databases like UniProt or DrugBank impossible without manual curation. Molecular identifiers receive no special treatment: SMILES strings, amino acid sequences, and CAS registry numbers are treated as text rather than validated against chemical and protein structure databases.

These limitations point to a fundamental gap between what similarity-based graph construction provides and what scientists require. Researchers need to trust their knowledge graphs. They need typed, directional, evidence-backed relationships grounded in established ontologies. They need to know that ``sotorasib \textsc{inhibits} KRAS G12C'' was \emph{understood} from the source text --- not inferred from vector proximity. The distinction between co-occurrence and comprehension is not merely academic; it determines whether a knowledge graph can support hypothesis generation, drug repurposing analysis, and regulatory evidence assembly, or whether it remains a visualization artifact.

The release of sift-kg \cite{sifkg2024}, a lightweight knowledge graph construction toolkit with built-in entity resolution, community detection, and interactive visualization, provided the missing infrastructure layer. Combined with Claude Code's agentic architecture --- which can dispatch parallel LLM subprocesses, enforce schemas via skill prompts, and invoke deterministic tools within a single execution flow --- the pieces were in place for a fundamentally different approach.

I present Epistract, a domain-specific agentic architecture that replaces similarity-based assembly with comprehension-based synthesis. In this system, LLM agents read scientific documents with domain understanding, extract entities and relations against a configurable schema grounded in over 40 biomedical ontologies, and validate molecular identifiers using deterministic tools including RDKit for chemical structures and Biopython for protein sequences. The resulting knowledge graphs capture meaning rather than co-occurrence. I validate this approach across six drug discovery domains --- neurogenetics, oncology, rare disease, immuno-oncology, cardiovascular inflammation, and GLP-1 receptor agonist competitive intelligence --- with zero retraining between domains, achieving 100\% entity type coverage and 93\% relation type coverage across all scenarios. The sixth scenario extends beyond PubMed abstracts to include patent documents sourced via Google Scholar and Google Patents (SerpAPI), demonstrating that the pipeline can extract molecular identifiers --- peptide sequences, CAS registry numbers, InChIKeys, and chemical formulas --- from patent filings. This multi-source workflow mirrors how a human researcher actually conducts competitive intelligence: searching multiple databases, cross-referencing findings, and building a structured knowledge base. The system is open-source and installable as a Claude Code plugin in a single command.

The remainder of this paper is organized as follows. Section~\ref{sec:architecture} describes the agentic architecture, extraction schema, and validation pipeline. Section~\ref{sec:schema} details the domain schema design and ontology grounding. Section~\ref{sec:evaluation} presents the evaluation methodology and results across the six therapeutic domains. Section~\ref{sec:comparison} compares the evolution from GraphRAG to Epistract, and Sections~6--8 discuss the human-AI collaboration process, availability, and future directions.
